\cotefile{156-3}
\scanpage{18}\mspage{17}
Sous tronçon \deleted{\illegible{1}}\inserted{fermé $I$} d'un tronçon $\unsure{L|\mathscr{L}}$ : sous ens. de la forme $[a,b]$ (avec $a<b$) ou $L_{\leq a}$ (avec $a$ pas le plus petit élément) ou $L_{\geq\deleted{\unsure{b}}a}$ (avec $b$ pas un plus grand élément), \inserted{ou $L$} : ce sont donc des \unsure{tronçons}, pour la structure de \unsure{biordre|biordonnés|bi-ordre} induite.
\inserted{On a $\partial I=\{a,b\}$ resp $\{a\}$ resp. $\{b\}$, resp $\emptyset$}
% slanted
\margin{Attention il s'agit de frontière \uline{relative} de $I$ dans $L$, et non du bord abs, \unsure{de|abs} $I$.}
\deleted{Donc \unsure{deux} Deux} Deux sous-tronçons fermés sont \deleted{\unsure{disjts|disj.}} en pos. rel. modérée si leurs frontières le sont, i.e. si \unsure{tt|les} \deleted{\unsure{él}} \unsure{point|lieu} frontière de l'un est \unsure{en} pos. rel. modérée à tt \unsure{lieu|point} frontière de l'autre.
\inserted{Sous-tronçon fermé large : soit un sous-tronçon, soit une partie uniponctuelle $\{a\}$, dans ce cas, sa frontière est $a$.}
\deleted{Alors leur int} \deleted{\unsure{est} \unsure{partie}, \unsure{relat.} modérée, \unsure{au cas où}} Deux tronçons fermés larges sont en pos. gén. \deleted{\unsure{sss|ssi}} leurs frontières \deleted{\unsure{donc son} s'il \unsure{le} est en \unsure{pos.} \unsure{gén.} modérée ou si} le sont.
% slanted
\margin{$(I\cap J)^{\unsure{\cdot|\prime|\circ}}\subset I\cup J^{\cdot}$ est vrai par argument topologique,}
Alors leur intersection est vide, \deleted{\unsure{ou un}} ou un \deleted{\illegible{1}} tronçon large, dont la frontière est contenue dans \deleted{\unsure{frontière} \unsure{large en pos. rel. modérée}} la réunion des frontières.
\deleted{\unsure{enfin|afin} \unsure{tt|les} pts frontière. \unsure{qui} \unsure{est} fermés deux tronçon fermé large : tout partie de la}

$\bigl[$\uline{Partie fermée} \uline{modérée} de $I$ : toute partie qui se met sous forme de réunion d'une famille finie de \unsure{ss-tronçons|ss.-tronçons} fermés larges mutuellement disjoints. et mut. \deleted{\unsure{en position}} rel. modérée.$\bigr]$

\unsure{Top.|Tp.} de $L$ \unsure{ici} : Prenant comme base d'ouverts les \unsure{«|``} $]a,b[$ pour $a<b$, et les $L_{>a}$ et $L_{<a}$ \unsure{(}. Alors \deleted{$\unsure{I}$} pour un sous-ps. tronçon \inserted{fermé}, $\partial\unsure{I|\widehat I}$ est la frontière
% slanted
\margin{\uline{NB} Il n'y a pas besoin d'introduire les $L_{>a}$ \unsure{et} si $L$ \unsure{n'a} pas de plus grand élément, et \unsure{itou|idem} pour $L_{<a}$ si \unsure{$a$|$L$} n'a pas de \unsure{plus} petit élément. \unsure{Dans ce cas|Dans} \illegible{2} $]a,b[$ \unsure{eux} \illegible{2} $]a,b[$.}
